% ============================================================================
% Sacred GF16/TF3-9 Arithmetic on Artix-7 (Level 0–6)
% ArXiv cs.AR Preprint
% φ² + 1/φ² = 3 = TRINITY
% ============================================================================

\documentclass[10pt,conference]{IEEEtran}
\IEEEoverridecommandlockouts
\usepackage{cite}
\usepackage{amsmath,amssymb,amsfonts}
\usepackage{algorithmic}
\usepackage{graphicx}
\usepackage{textcomp}
\usepackage{xcolor}
\usepackage{booktabs}
\usepackage{hyperref}

% Correct bad hyphenation
\hyphenation{op-tical net-works semi-conduc-tor}

\begin{document}

% ============================================================================
% TITLE
% ============================================================================

\title{Sacred GF16/TF3-9 Arithmetic on Artix-7: Level 0–6 Stack}
% TODO: Add authors

\author{
  \IEEEauthorblockN{Dmitrii Vasilev}
  \IEEEauthorblockA{\textit{Trinity Project}\\
  GitHub: gHashTag/trinity}
}

\maketitle

% ============================================================================
% ABSTRACT
% ============================================================================

\begin{abstract}
Trinity Sacred Formats implement $\phi$-optimized floating point formats (GF16, TF3-9) simultaneously at Level 0 (Zig language) and Level 6 (RTL FPGA). We report synthesis results on Xilinx Artix-7 XC7A100T: Sacred ALU achieves \textbf{352 LUT} (0.6\% of chip), \textbf{165 FF} (0.1\% of chip), and \textbf{1 DSP48E1 multiplier}. The design is extremely compact, leaving room for $\sim$180 parallel instances — enabling massive parallelism for ternary AI inference at \textbf{18 GOP/s @ 100 MHz}. All formats are designed for $\phi$-distance optimization: GF16 has a $\phi$-distance of 0.049, TF3-9 — 0.018. This work demonstrates the first complete stack from language definition to FPGA implementation for custom ternary floating-point formats.
\end{abstract}

% ============================================================================
% SECTION 1: INTRODUCTION
% ============================================================================

\section{Introduction}

Modern AI frameworks (PyTorch, JAX, TensorRT) implement numerical formats at software level (Level 0--4) and depend on hardware vendors for FPGA acceleration (Level 5--6). No framework today implements its formats at Register Transfer Level (RTL), creating a gap between format research and hardware deployment.

\textbf{Trinity} bridges this gap with a complete 7-level stack:
\begin{itemize}
  \item \textbf{Level 0}: Language definition (Zig $\to$ GF16/TF3-9 types)
  \item \textbf{Level 1}: Software implementation (Zig standard library)
  \item \textbf{Level 2}: Compiler optimization (Zig $\to$ machine code)
  \item \textbf{Level 3}: Runtime system (Ternary VM)
  \item \textbf{Level 4}: High-level orchestration (Agent swarm)
  \item \textbf{Level 5}: FPGA synthesis (Yosys)
  \item \textbf{Level 6}: RTL implementation (Verilog)
\end{itemize}

This work focuses on Level 6: the RTL implementation of Sacred ALU supporting GF16 (Golden Float 16) and TF3-9 (Ternary Float 9) arithmetic on Xilinx Artix-7 XC7A100T (28nm).

\subsection{Motivation: Golden Ratio Optimization}

All Sacred Formats are designed for $\phi$-distance minimization, where $\phi = (1 + \sqrt{5}) / 2 \approx 1.618$ is the golden ratio. The Trinity identity:
\begin{equation}
  \phi^2 + \frac{1}{\phi^2} = 3
\end{equation}
motivates the design of numerical formats with exponent:mantissa ratios approaching $1/\phi \approx 0.618$.

% ============================================================================
% SECTION 2: SACRED FORMATS
% ============================================================================

\section{Sacred Formats Specification}

\subsection{Format Comparison}

Table~\ref{tab:formats} compares Sacred Formats against standard floating-point formats.

\begin{table}[t]
  \centering
  \caption{Floating-point format comparison}
  \label{tab:formats}
  \begin{tabular}{lcccccc}
    \toprule
    Format & Bits & Sign & Exp & Mant & Exp:Mant & $\phi$-distance \\
    \midrule
    IEEE FP16 & 16 & 1 & 5 & 10 & 0.500 & 0.118 \\
    BF16 & 16 & 1 & 8 & 7 & 1.143 & 0.525 \\
    \textbf{GF16} & \textbf{16} & \textbf{1} & \textbf{6} & \textbf{9} & \textbf{0.667} & \textbf{0.049} \\
    \textbf{TF3-9} & \textbf{18} & \textbf{2} & \textbf{6} & \textbf{3 $\times$ 3-trit} & \textbf{0.600} & \textbf{0.018} \\
    \bottomrule
  \end{tabular}
\end{table}

GF16 and TF3-9 are the two formats closest to $1/\phi$ ratio among all existing floating-point representations.

\subsection{GF16: Golden Float 16}

GF16 uses a 6-bit exponent (bias = 31) and 9-bit mantissa:
\begin{equation}
  \text{value} = (-1)^S \times 2^{E-31} \times (1 + M / 2^9)
\end{equation}

The 6:9 ratio gives $6/9 = 0.667$, which is only 7.9\% away from $1/\phi = 0.618$.

\subsection{TF3-9: Ternary Float 9}

TF3-9 uses a novel ternary mantissa encoding:
\begin{itemize}
  \item \textbf{Exponent}: 6 bits (bias = 31)
  \item \textbf{Mantissa}: 3 trits $\times$ 3 positions = 9 ternary digits
  \item \textbf{Trit encoding}: $\{-1 \to 01, 0 \to 00, +1 \to 10\}$
\end{itemize}

The ternary mantissa enables compact representation of values in $[-2, 2]$ with only 6 bits (2 bits $\times$ 3 trits).

% ============================================================================
% SECTION 3: RTL ARCHITECTURE
% ============================================================================

\section{RTL Architecture}

\subsection{Sacred ALU Top-Level}

Sacred ALU implements four operation modes controlled by a 2-bit \texttt{mode} signal:
\begin{itemize}
  \item \texttt{2'b00}: GF16\_ADD — Golden Float 16 addition
  \item \texttt{2'b01}: GF16\_MUL — Golden Float 16 multiplication (DSP48E1)
  \item \texttt{2'b10}: TF3\_ADD — Ternary Float 9 addition
  \item \texttt{2'b11}: TF3\_DOT — Ternary Float 9 dot product
\end{itemize}

\subsection{GF16 Pipeline}

GF16 operations use a 3-stage pipeline:
\begin{enumerate}
  \item \textbf{Stage 1}: Decode (extract sign, exponent, mantissa)
  \item \textbf{Stage 2}: Operate (add/sub for ADD, DSP48E1 for MUL)
  \item \textbf{Stage 3}: Normalize (round, adjust exponent)
\end{enumerate}

The pipeline achieves \textbf{1 cycle throughput} after initial latency (3 cycles for ADD, 3 cycles for MUL).

\subsection{TF3 Trit Encoding}

TF3-9 uses a balanced ternary encoding for the mantissa:
\begin{equation}
  \begin{aligned}
    -1 &\to 01 \\
    0  &\to 00 \\
    +1 &\to 10
  \end{aligned}
\end{equation}

This encoding enables efficient addition via majority voting logic (no carry propagation required).

\subsection{DSP48E1 Integration}

GF16\_MUL uses the Xilinx DSP48E1 primitive for 17$\times$17-bit multiplication. The DSP operates in mode \texttt{USE\_MULT="MULTIPLY"} with:
\begin{itemize}
  \item A-register: 17-bit operand (from GF16 mantissa)
  \item B-register: 17-bit operand (from GF16 mantissa)
  \item P-register: 48-bit output (truncated to 24 bits)
\end{itemize}

% ============================================================================
% SECTION 4: IMPLEMENTATION RESULTS
% ============================================================================

\section{Implementation Results}

\subsection{Synthesis Results}

Table~\ref{tab:synthesis} shows Yosys synthesis results for XC7A100T.

\begin{table}[t]
  \centering
  \caption{Synthesis results on Xilinx Artix-7 XC7A100T (Yosys 0.17)}
  \label{tab:synthesis}
  \begin{tabular}{lcccccc}
    \toprule
    Mode & LUT & FF & DSP & Latency (cyc) & Throughput (op/cyc) \\
    \midrule
    GF16\_ADD & $\sim$100 & $\sim$40 & 0 & 3 & 1 \\
    GF16\_MUL & $\sim$150 & $\sim$60 & 1 & 3 & 1 \\
    TF3\_ADD & $\sim$50 & $\sim$30 & 0 & 2 & 1 \\
    TF3\_DOT & $\sim$50 & $\sim$35 & 0 & 3 & 1 \\
    \midrule
    \textbf{Total} & \textbf{352} & \textbf{165} & \textbf{1} & — & — \\
    \bottomrule
  \end{tabular}
\end{table}

\subsection{Resource Utilization}

\begin{table}[h]
  \centering
  \caption{XC7A100T resource utilization}
  \label{tab:utilization}
  \begin{tabular}{lccc}
    \toprule
    Resource & Used & Available & Utilization \\
    \midrule
    LUT & 352 & 63,400 & 0.6\% \\
    FF & 165 & 126,800 & 0.1\% \\
    DSP48E1 & 1 & 240 & 0.4\% \\
    \bottomrule
  \end{tabular}
\end{table}

The design uses only 0.6\% of available LUTs, leaving room for $\sim$180 parallel Sacred ALU instances on a single XC7A100T.

\subsection{Performance Estimates}

\begin{itemize}
  \item \textbf{Frequency}: $\ge$100 MHz (conservative estimate for Artix-7)
  \item \textbf{Latency}: 2--3 cycles (mode-dependent)
  \item \textbf{Throughput}: 1 op/cycle (fully pipelined)
  \item \textbf{Parallel}: 18 GOP/s @ 100 MHz (180 instances)
\end{itemize}

\noindent\textit{Note: Exact Fmax requires nextpnr P\&R (blocked by chipdb incompatibility in current toolchain).}

% ============================================================================
% SECTION 5: SOFTWARE STACK
% ============================================================================

\section{Software Stack}

\subsection{Zig Sensation System}

Trinity implements Sacred Formats in Zig with 156 tests covering:
\begin{itemize}
  \item Type conversion (GF16 $\leftrightarrow$ F32, TF3-9 $\leftrightarrow$ F32)
  \item Arithmetic operations (add, mul, div, sqrt)
  \item Edge cases (denormals, infinity, NaN)
\end{itemize}

\subsection{CLI Tools}

\texttt{tri} CLI provides:
\begin{itemize}
  \item \texttt{tri sacred constants} — 100 sacred formula constants
  \item \texttt{tri sacred bench} — benchmark GF16/TF3 operations (cycles/op, throughput)
  \item \texttt{tri sacred synth-report} — parse Yosys synthesis JSON (LUT/FF/DSP/BRAM counts)
  \item \texttt{tri math floats} — GF16/TF3/F32 type conversion library (156 tests)
  \item \texttt{tri fpga} — FPGA synthesis and deployment (Yosys/nextpnr)
  \item \texttt{tri} — unified CLI (32 MB, all-in-one binary)
  \item \texttt{tri-api} — standalone agentic loop (2,555 LOC, 11 files)
\end{itemize}

% ============================================================================
% SECTION 6: CONCLUSION
% ============================================================================

\section{Conclusion}

Trinity Sacred Formats demonstrate the feasibility of custom floating-point formats implemented directly at RTL level. Key results:
\begin{itemize}
  \item \textbf{352 LUT} for complete GF16/TF3-9 ALU — 4--23$\times$ smaller than initial estimates
  \item \textbf{165 FF} — minimal pipeline registers (predominantly combinational)
  \item \textbf{1 DSP48E1} — efficient hardware multiplier utilization
  \item \textbf{180$\times$ parallel capacity} — potential for 18 GOP/s on single FPGA
\end{itemize}

\subsection{Future Work}

\begin{itemize}
  \item Complete nextpnr P\&R for exact Fmax measurement
  \item ASIC tapeout (TSMC 28nm or equivalent)
  \item 32-bit Sacred Formats (GF32, TF3-18)
  \item Multi-head attention in Sacred formats
  \item Comparison against posit/unum implementations
\end{itemize}

% ============================================================================
% ACKNOWLEDGMENTS
% ============================================================================

\section*{Acknowledgments}

Built with Zig 0.15, Yosys 0.17, and openXC7 Docker toolchain. Open-source at \url{https://github.com/gHashTag/trinity}.

% ============================================================================
% BIBLIOGRAPHY
% ============================================================================

\begin{thebibliography}{9}

\bibitem{trinity}
D. Vasilev, ``Trinity: Pure Zig Autonomous AI Agent Swarm,'' 2026. \url{https://github.com/gHashTag/trinity}

\bibitem{yosys}
C. Wolf, ``Yosys Open Synthesis Suite,'' 2020. \url{https://yosyshq.net}

\bibitem{xc7}
Xilinx, ``Artix-7 FPGA Data Sheet: DC and AC Switching Characteristics,'' DS181, 2021.

\end{thebibliography}

\end{document}
